\begin{proposition}[C.3.21 (Proved): the actual delta-line square]
For the actual four-manifold and line bundle in C.3.20,
\[
 s_\delta=-2.
\]
The comparison is with the actual relative Picard section of
the circle bundle; it does not assume that the compact-core
model of Part B fills the original six-sphere.
\end{proposition}
\begin{proof}
We first identify the base precisely. On the Milnor page used in
C.3.18, choose its phase so that
$F_\epsilon=(1-\epsilon)z_2^3-z_1^4$ is positive real.
Positive weighted rescaling identifies its interior with
\[
 C^\circ:\quad b^3=a^4+1,
 \qquad
 a=\frac{z_1}{F_\epsilon^{1/4}},\qquad
 b=\frac{(1-\epsilon)^{1/3}z_2}{F_\epsilon^{1/3}}.
\]
Conversely, every point of this affine curve has a unique
positive weighted rescaling to the unit three-sphere. A different
page phase changes this by constant complex factors and gives
the same marked cover. The compactification $C$ has genus three
and one point at infinity. Its actual map to the original
$t$-line is
\[
 t_\epsilon=\frac{a^4+1}{1+\epsilon a^4}.
\]
It is ramified with orders three and four over $0$ and $1$,
and with order twelve at the added point over $t=1/\epsilon$.
The twelve points over $t=\infty$ satisfy
$a^4=-1/\epsilon$ and are unramified. Thus this is the
actual page cover; the cover ramified twelvefold over the
original cusp is its $\epsilon=0$ comparator, not the same
marked map when $\epsilon>0$.

Keep the product modification inside $|q|<r_0<\epsilon$.
The page equation there has twelve disjoint graph discs
$\theta=\theta_j(q)$, since its derivative in $\theta$ is
twelve. Consequently $M$ consists of twelve nodal elliptic
neighbourhoods attached to an entirely untouched raw torus
quotient. Denote the latter part by $M^\circ$ and its base by
$C_0$; it is the page with these twelve discs removed. It has
nonempty boundary and retracts to a graph.

We identify the actual circle bundle on this unmodified part.
On an ordinary marked chart let $g$ be a real lift of its
$\gamma$ coefficient and set
\[
 \xi=\zeta_1-6\mu g,\qquad
 \eta=\zeta_2-\beta g.
\]
Then $\xi=\tau u+w$ and $\eta=\mu u+d$. Exponentiating
$\eta$ gives the actual degree-zero line on the first elliptic
torus with multipliers
\[
 (\xi,z)\sim(\xi+\tau,e^{2\pi i\mu}z),\qquad
 (\xi,z)\sim(\xi+1,z).
\]
These are the Poincar\'e multipliers of the original normal
function $P$, up to the single dualization convention for the
associated circle line. The identification is an identification
of the varying fibrewise Picard section, not just its fibre
degree. Under a change of marked elliptic periods the maps are
fibrewise affine holomorphic and the line maps are its
semiabelian lifts. The translations may depend smoothly on
$g$; no global holomorphic-torsor assertion is needed.

Over the graph base $C_0$, the affine elliptic torsor has a
smooth section. Indeed its obstruction lies in the twisted
group $H^2(C_0,\Lambda)$, which is zero; equivalently use the
sequence from the lattice local system to its real vector
bundle and then to the torus bundle, whose vector-bundle
section sheaf is fine. Identify the torsor with its Jacobian
using such a section. Translation of a degree-zero Poincar\'e
line changes it by a line pulled back from the base. Two
fibrewise holomorphic lines with the same varying Picard
section likewise differ by a base line: their quotient has a
one-dimensional space of fibrewise holomorphic sections and
evaluation trivializes it fibrewise. Since $H^2(C_0;\Z)=0$,
this base line has zero topological Chern class. The actual
class on $M^\circ$ is therefore the Poincar\'e class of
$P\circ t_\epsilon$, up to sign, including its mixed
base--fibre component.

There is no hidden correction supported at the twelve modified
nodes. By excision and Poincar\'e--Lefschetz duality, each
local summand of
\[
 H^2(M,M^\circ;\Z)
 \simeq\bigoplus_{j=1}^{12}H^2(E_{I_1},\partial E_{I_1};\Z)
\]
is $\Z$, generated by the Thom class of a regular elliptic
fibre. The regular fibre is primitive in $H_2(E_{I_1};\Z)$,
as its degree-one specialization generates the nodal fibre's
two-cycle. Its image in $H^2(M;\Z)$ is the pullback of the
point class of the base, and is zero because that base has
boundary. The restriction $H^2(M;\Z)\to H^2(M^\circ;\Z)$
is consequently injective. This addresses actual local
clutching classes, rather than only their holonomy.

Trim a smaller boundary collar, which does not change the
relative square because its line is trivial, and identify the
remaining annular end of the page with the punctured disc at
the added point of $C$. Fill it by its regular marked
$T^2\times D^2$. The boundary affine torus bundle is
fibrewise complex-trivialized by a path of translations.
After matching the Picard section $P$, a remaining line
identification is a scalar base-line gauge. Thus one can choose
the product-cap Poincar\'e extension whose section is
$P\circ t_\epsilon$ near the added point; the displayed
multipliers are regular there. The line is topologically
trivial on the cap, and its square contribution is zero.
Boundary gauge changes cannot change the relative square, by
C.3.20. Call the resulting
closed smooth genus-one Lefschetz surface $\widehat M$ and its
line $\widehat{\mathcal L}$. Its square is the same relative
square $s_\delta$; the cap contributes zero.

We spell out why the square can now be compared with the
Jacobian without a projectivity or multisection assumption on
$\widehat M$. Let $j:C^*=C\setminus\{12\text{ nodal values}\}
\hookrightarrow C$ and let $\Lambda=R^1\pi_*\Z|_{C^*}$.
For a nodal elliptic fibre, degree-one specialization identifies
$R^2\pi_*\Z$ with the constant sheaf $\Z$, while degree one
is the invariant subgroup of the unipotent monodromy. Thus
\[
 R^1\pi_*\Z=j_*\Lambda,\qquad R^2\pi_*\Z=\Z.
\]
These sheaves and their fibre intersection pairing agree for
the actual surface and its Jacobian. The degree-zero class
$c_1(\widehat{\mathcal L})$ lies in the first Leray filtration
step. Its component in
$H^1(C,j_*\Lambda)$ is the connecting class of the varying
Picard section $\pm P$. This holds on $C^*$ by the preceding
actual multiplier calculation. It holds on all of $C$ because
\[
 H^1(C,j_*\Lambda)\longrightarrow H^1(C^*,\Lambda)
\]
is injective, by the first terms of the Leray sequence for $j$.
Local nodal choices can therefore add only a base-filtration
class; the excision argument gives the same conclusion directly.

Multiplicativity of Leray now determines the square. In degree
four its only nonzero graded piece is
$H^2(C,R^2\pi_*\Z)=\Z$. The square of a class in filtration
one is the cup square of its $H^1(C,j_*\Lambda)$ component,
using the fibre intersection pairing. Adding filtration-two
base classes contributes only filtration three or four, both
zero on this base. Hence the actual square is exactly the
square of the Poincar\'e class on the Jacobian; dualizing the
line does not change it.

Finally let $J_\epsilon\to C$ be the relatively minimal
Jacobian obtained by pulling back the original elliptic surface
along $t_\epsilon$. It has the explicit affine model
\[
 Y^2=X^3-3\lambda bX+2\lambda^2a^2,
 \qquad P=(0,-\sqrt2\lambda a),\qquad
 \lambda=1-\epsilon.
\]
Writing $h=1+\epsilon a^4$, the substitution from the original
Weierstrass coordinates is
$x=(ab^2/h)^2X$, $y=(ab^2/h)^3Y$.
The discriminant is
\[
 \Delta=1728\lambda^3(1+\epsilon a^4).
\]
It has twelve simple zeroes on $C^\circ$, and $c_4$ is a unit
at all of them. They are precisely the twelve $I_1$ fibres.
At the unique point at infinity, $a,b$ have pole orders
three and four. Scaling $X,Y$ by a local parameter to orders
two and three gives a regular Weierstrass model whose
discriminant is a unit when $\epsilon\ne0$; the point at
infinity is therefore a good fibre. This also shows that the
fundamental line has degree one, so
$\chi(\mathcal O_{J_\epsilon})=1$.

The displayed section $P$ is finite on every affine chart and
also in this integral chart at infinity, hence it is disjoint
from the zero section. At a simple discriminant zero it is a
smooth point of the nodal fibre: a node has $Y=0$, whereas
$a\ne0$ there and $Y(P)=-\sqrt2\lambda a\ne0$.
All singular fibres are irreducible, so there is no component
correction. These explicit identities are independently checked
in \texttt{ab\_push\_actual\_line\_comparison.py}.

The divisor $P-O$ has fibre degree zero and represents the
required relative Picard section, up to a base twist. Both
sections have square $-\chi=-1$ and $P\cdot O=0$. Therefore
\[
 (P-O)^2=-1-1-2(P\cdot O)=-2.
\]
Base twists do not change this square. The actual Leray
comparison above proves $s_\delta=-2$, as asserted.
\end{proof}
